\documentclass[12pt]{article}
\usepackage{fontspec}
\usepackage{graphicx}
\usepackage{pst-plot}
\usepackage{scrextend}
\usepackage{advdate}
\usepackage{listings}
\usepackage{tikz}
\usepackage{amssymb}
\usepackage{setspace}
\usepackage{color}
\usepackage{soul}
\usepackage{caption}
\usepackage{fancyhdr}
\usepackage{pgfornament}
\usepackage[many]{tcolorbox}
\usepackage{collectbox}
\usepackage{mdframed}
\usepackage{xcolor}
\usepackage{mathtools}
\usepackage{wasysym}
\usepackage[hidelinks]{hyperref}
\usepackage[final]{pdfpages}
\usepackage[left=0.5in,right=0.5in,top=0.5in,bottom=0.5in,footskip=15mm]{geometry}

\newfontfamily{\ubuntu}[Path=./fonts/,]{UbuntuMono-Regular}
\newfontfamily{\ubuntui}[Path=./fonts/,]{UbuntuMono-Italic}
\newfontfamily{\ubuntub}[Path=./fonts/,]{UbuntuMono-Bold}

\definecolor{blue1}{RGB}{115, 3, 252}
\definecolor{blue2}{RGB}{3, 53, 252}
\definecolor{orange1}{RGB}{255, 68, 0}
\setmainfont{Arial}
\onehalfspacing 
\begin{document}
\pagenumbering{gobble}

\hfill \textbf{Rocco DeVito}

\hfill \today

\vspace{2mm}

\noindent {\Large \textbf{\textcolor{blue1}{The Polynomial Ring $\mathbb{Z}[x]$} (Part One)  \hrulefill}}

\vspace{8mm}

\noindent A term that you should become familiar with is \textbf{ostensive definition}. To define a word ostensively, you give examples of things that would be classified as whatever the word is. There are also \textbf{formal definitions} in mathematics, but they are often extremely difficult to comprehend without sufficient experience. For this reason, I will resort to using ostensive definitions as much as possible. Note, however, that deep mathematics requires formal definitions. Ostensive definitions are only for pedagogical purposes.

\vspace{5mm}

\noindent Let me ostensively define what \textbf{polynomials} are. They are things like these:
\begin{align*}
&x^2 + 2x + 3, \\
&x^4-6x^3+9x-1, \\
&x^{11}-5x, \\
&x^{27}-8x^{14}+9x^3-x
\end{align*}
\noindent The numbers in front of the powers of $x$ are called \textbf{coefficients}. For example, the coefficient on the $x^3$ term in the second polynomial above is $-6$. Knowing this word is crucial for fluently conversing in the language of mathematics. Do not worry about what $x$ is. It is unknown, and we are \underline{\textbf{not}} trying to solve for it.

\vspace{5mm}

\noindent Take note of the following: Polynomials themselves are \underline{\textbf{not}} equations. There is such a thing as a polynomial equation, but the entities above have no equals signs in them. Thus, they are not equations. To further clarify what polynomials are, let me provide you with some \textbf{non-polynomials} (things that aren't polynomials):
\begin{align*}
&e^{2x}+7, \\
&\ln(5x+2)+\sqrt{x},\\
&e^x + x^6+2x-1,\\
&6\cos(x)+ 8\sin(x)
\end{align*}

\noindent If you take all the polynomials with coefficients in $\mathbb{Z}$, you get a ring that we denote by $\mathbb{Z}[x].$ You can call it \textbf{the ring of polynomials with coefficients in $\mathbb{Z}$}. Remember: a ring is a set equipped with \underline{\textbf{two}} operations. For today's lesson, we will only concern ourselves with the addition operation. We will save multiplication for next time. Let me ostensively define polynomial addition. Let's add together the polynomials $4x^3-5x^2+x-7$ and $-9x^7+5x^3+12x^2+9$. I'm going to start by putting each of these in parentheses for the sake of readability and distinction. We have
\begin{align*}
\textcolor{red}{(}4x^3-5x^2+x-7\textcolor{red}{)} + \textcolor{red}{(}-9x^7+5x^3+12x^2+9\textcolor{red}{)} &= -9x^7 + 4x^3 + 5x^3 - 5x^2 + 12x^2 + x - 7 + 9 \\
&= -9x^7 + (4x^3 +5x^3) + (-5x^2 + 12x^2) + x + (-7 + 9) \\
&= -9x^7 + 9x^3 + 7x^2 + x + 2.
\end{align*}
\noindent Look at the above calculation very carefully. To perform polynomial addition, you simply add coefficients on corresponding powers of $x$ together. If one of your polynomials had $6x^3$ in it and the the other had $7x^3$ in it, your resulting polynomial would have $13x^3$ in it. Note that terms like $3x^8$ and $2x^2$ cannot be added together, because their powers of $x$ are not the same. 

\vspace{5mm}

\begin{center}
{\large\textbf{Exercises}}
\end{center}

\vspace{3mm}

\noindent \textcolor{blue1}{(1.)} Find the sum of $x^4+9x^2$ and $x^5 - 2x^4 + 2x^2$.

\vspace{30mm}

\noindent \textcolor{blue1}{(2.)} Find the difference of $2x^3 + 4x^2$ and $7x^3 + 2x^2$.

\vspace{30mm}

\noindent \textcolor{blue1}{(3.)} Prove (informally) that $\mathbb{Z}[x]$ forms an Abelian group under polynomial addition (``Abelian'' means ``commutative'').





\end{document}
